% Project report template for CS6370 — Parts 1–4 (Cranfield IR)
\documentclass[11pt,a4paper]{article}
\usepackage[utf8]{inputenc}
\usepackage[T1]{fontenc}
\usepackage{microtype}
\usepackage{geometry}
\geometry{margin=1in}
\usepackage{graphicx}
\usepackage{amsmath,amssymb}
\usepackage{booktabs}
\usepackage{hyperref}
\usepackage{enumitem}

\title{Project Report: Improving a Vector-Space IR System\\(Parts 1--4)}
\author{Name: \underline{\hspace{6cm}} \\
Roll No.: \underline{\hspace{4cm}}}
\date{}

\begin{document}
\maketitle
\begin{center}
\textbf{Course:} CS6370 -- Natural Language Processing\\
\textbf{Dataset:} Cranfield (1400 documents, 225 queries)
\end{center}

\begin{abstract}
This document is a working LaTeX template for the CS6370 project report. It contains
the report structure required by the assignment and an initial draft for Parts 1--4.
Fill in code snippets, figures, tables, and results produced by your implementation.
\end{abstract}

\tableofcontents
\newpage

\section{Introduction}
This project implements and analyses an Information Retrieval (IR) system based on the
Vector Space Model (VSM) using the Cranfield dataset. Parts 1--4 of the project cover:
preprocessing and indexing, the TF--IDF representation and retrieval, evaluation of the
basic IR system using standard metrics, and an analysis of VSM limitations with concrete
examples. Subsequent parts will investigate improvements to the baseline system.

\subsection{Objectives}
\begin{itemize}[noitemsep]
  \item Build a reproducible pipeline for preprocessing (segmentation, tokenization,
    stopword removal, stemming/lemmatization) and inverted index construction.
  \item Represent documents and queries with TF--IDF vectors and implement retrieval
    using cosine similarity and ranking.
  \item Evaluate retrieval with Precision@k, Recall@k, F0.5@k, AP@k, nDCG@k, MAP and MRR.
  \item Analyse limitations of VSM and propose concrete improvements for the project.
\end{itemize}

\section{Limitations of the Vector Space Model (VSM)}
Below we summarise key limitations of the classical TF--IDF VSM and illustrate them
using examples drawn from the Cranfield dataset and the assignment's toy examples.

\subsection{Term independence and lack of semantics}
VSM treats terms as independent dimensions. This prevents the model from capturing
semantic relations (synonymy) or disambiguating multiple senses (polysemy). For
example, the word \emph{"star"} is ambiguous: it can mean an astronomical star or
a movie star. A query like ``movie star'' should ideally retrieve documents about
celebrities (assignment's d2), but a naive VSM may return documents about the sun
if those documents share the token \textit{star} and related words.

\subsection{Zero-result / OOV (Out-of-Vocabulary) problem}
When a query contains terms not present in the indexed vocabulary, the standard
TF--IDF implementation yields zero overlap and hence no matches. For example, a
specialised technical term absent from the Cranfield vocabulary will cause the
system to return no results unless query-expansion or robust matching is used.

\subsection{Sensitivity to surface form and morphology}
Differences such as plurals, inflections, or spelling variants (``computer'' vs
``computers'' or misspellings) reduce matching quality unless normalization like
stemming or lemmatization is applied. Stemming may help but can also over-conflate
distinct terms.

\subsection{Document length and term weighting issues}
Raw TF favors longer documents; TF--IDF mitigates this but requires careful
normalization (e.g., cosine normalization). Ranking can still be suboptimal when
documents differ widely in length or in the amount of topical vs filler text.

\subsection{Lack of phrase/positional information}
VSM typically discards term order and proximity, making it unable to distinguish
``information retrieval'' from documents that contain the words far apart.

\section{Proposed Approaches to Address These Issues}
This section lists candidate methods to address the limitations above. These are
not all required to be implemented; they serve as a roadmap for Part 5 (project).

\subsection{Semantic representations}
\begin{itemize}[noitemsep]
  \item Latent Semantic Analysis (LSA): SVD on the term-document matrix to capture
    latent concepts (reduces synonymy effects).
  \item Word embeddings (Word2Vec, GloVe) or contextual embeddings (BERT): represent
    terms and queries in dense semantic space and use cosine similarity there.
\end{itemize}

\subsection{Improved retrieval models}
\begin{itemize}[noitemsep]
  \item BM25 / Language models: probabilistic scoring functions that often outperform
    plain TF--IDF in practice.
  \item Query expansion and pseudo-relevance feedback: expand queries with related
    terms or top-retrieved terms to improve recall.
\end{itemize}

\subsection{Robust preprocessing}
Better tokenization, lemmatization, spelling correction, and handling of multi-word
expressions (n-grams / phrases) reduce surface-form mismatches.

\section{Dataset Description}
We use the Cranfield dataset provided with the assignment. Key statistics (computed
from the provided JSON files):
\begin{itemize}[noitemsep]
  \item Documents: 1400 (see cranfield/cran_docs.json)
  \item Queries: 225 (see cranfield/cran_queries.json)
  \item Relevance judgments: cranfield/cran_qrels.json (use relevance > 0 as relevant)
\end{itemize}

Include in this section a short description of the JSON fields you use (e.g., id,
title, body) and any preprocessing steps (lowercasing, punctuation removal, tokenization,
stopword list, stemmer choice).

\subsection{Practical notes}
Place code and outputs in a single folder for submission: \texttt{RollNo.zip} containing
the code folder and a PDF with textual answers, as required by the assignment.

\section*{How to compile this template}
Run the following from the folder containing \texttt{project_report.tex}:
\begin{verbatim}
pdflatex project_report.tex
pdflatex project_report.tex
\end{verbatim}

\section*{Next steps}
Complete Part 1--4 artifacts by:
\begin{itemize}[noitemsep]
  \item Adding the implemented code screenshots and key snippets into an \texttt{appendix}.
  \item Populating the evaluation section with plots for metrics vs k (Part 3).
  \item Adding concrete retrieval examples (tables) that show failure cases for VSM.
\end{itemize}

\bibliographystyle{plain}
\begin{thebibliography}{9}
\bibitem{salton1996}
G. Salton and C. Buckley. Term-weighting approaches in automatic text retrieval.
\emph{Information Processing & Management}, 1996.
\end{thebibliography}

\end{document}
